\bookchapter{Capítulo 69: Manual de consulta rápida y recursos}

\chapterintro{Este capítulo es un manual de consulta rápida que puede usarse por sí solo. Cuando no recuerdes cómo se escribe algo o necesites decidir rápido, ven directo aquí. Trátalo como una «tarjeta de bolsillo»: código de referencia, preguntas frecuentes y recursos de aprendizaje, sin saltarte una página.}

\booksection{Tarjeta de referencia rápida de código}

\begin{center}
\begin{tabularx}{\linewidth}{|X|X|}
\hline
\textbf{Qué hacer} & \textbf{Código clave} \tabularnewline
\hline
Estado privado de la página & @State private var x = valorInicial \tabularnewline
Lógica entre páginas & \texttt{@Observable class VM \{ var x = 0 \}} \tabularnewline
Maquetación vertical/horizontal & \texttt{VStack \{ \}}/\texttt{HStack \{ \}} \tabularnewline
Espacio elástico & \texttt{Spacer()} \tabularnewline
Lista & List \{ ForEach(array) \{ ... \} \} \tabularnewline
Navegar entre páginas & \texttt{NavigationStack + NavigationLink} \tabularnewline
Ventana inferior & .sheet(isPresented: \$show) \{ contenido \} \tabularnewline
Cerrar la página actual & \texttt{@Environment(\textbackslash{}.\allowbreak dismiss) var dismiss} \tabularnewline
Traer datos de la red & \texttt{try await URLSession.\allowbreak shared.\allowbreak data(from: url)} \tabularnewline
Interpretar JSON & \texttt{struct X: Codable}+\texttt{JSONDecoder().\allowbreak decode} \tabularnewline
Guardar configuraciones pequeñas & UserDefaults.standard.set(valor, forKey: "clave") \tabularnewline
Guardar datos estructurados & \texttt{@Model}+\texttt{modelContainer}+\texttt{@Query} \tabularnewline
Adaptación al modo oscuro & Usa colores del sistema (.primary / Color(.systemGray6)), no colores fijos \tabularnewline
Ejecutar pruebas & CmdU \tabularnewline
\hline
\end{tabularx}
\end{center}

\begin{center}\small\color{gray}Tarjeta de referencia rápida de desarrollo iOS (edición de agosto de 2026)\end{center}

\booksection{Preguntas frecuentes (FAQ)}

P1: No tengo Mac, ¿puedo aprender desarrollo de iOS?

Puedes aprender el lenguaje, pero no hacer desarrollo completo. Con Swift Playgrounds en el iPad/iPhone puedes terminar todo el contenido del capítulo 61; pero Xcode es exclusivo de macOS y al final necesitas un Mac. Si el presupuesto es corto, puedes alquilar primero un Mac en la nube o comprar un Mac mini de segunda mano como transición, y decidir cuánto invertir cuando domines el lenguaje.

P2: ¿Cuánto cuesta aprender desarrollo de iOS?

Fase de aprendizaje: 0 (Xcode es gratuito y una cuenta Apple gratuita permite depurar en un dispositivo real). Fase de publicación: la cuenta de desarrollador cuesta 99 dólares al año (unos 700 yuanes). Sumando la inversión en un Mac, ese es todo el coste. Frente a lo que se recibe a cambio —en el App Store, hasta una pequeña utilidad encuentra usuarios que pagan—, la cuenta sale a favor.

P3: ¿Qué conviene aprender primero, SwiftUI o UIKit?

SwiftUI. En 2026 los proyectos, tutoriales y ofertas de empleo nuevos giran en torno a SwiftUI, y este libro lo usa de principio a fin. UIKit es el framework de la generación anterior y sigue vivo en muchísimos proyectos antiguos: cuando ya domines SwiftUI, añade UIKit según lo necesites (Apple ofrece tutoriales de interoperabilidad) y te pondrás rápido. No inviertas el orden.

P4: ¿Es fácil encontrar trabajo en desarrollo iOS?

El panorama de 2026: los puestos básicos de «solo saber dibujar interfaces» tienen mucha competencia, pero los desarrolladores que «dominan Swift + SwiftUI y tienen apps publicadas» siguen siendo muy buscados. El salario medio de los puestos iOS se mantiene año tras año a la cabeza entre las distintas tecnologías. Mientras estás aprendiendo, tu mejor currículum es esa app hecha por ti que ya está en la tienda.

P5: ¿Qué hago si el proceso de revisión rechaza mi app?

No entres en pánico: es totalmente normal —la probabilidad de que el primer envío sea rechazado no es baja. Apple suele escribir el motivo del rechazo con mucho detalle (qué pantalla, qué directriz): corrige punto por punto, responde con una explicación en App Store Connect y vuelve a enviar. Las tres causas de rechazo más habituales son: no explicar con claridad para qué se pide cada permiso, bugs en la interfaz y metadatos (capturas/descripción) que no coinciden con la app real. Si revisas estos tres puntos por adelantado, te ahorras una ronda de espera.

\booksection{Atajos de teclado de Xcode: guía de referencia rápida}

\begin{center}
\begin{tabularx}{\linewidth}{|X|X|}
\hline
\textbf{Acción} & \textbf{Atajo} \tabularnewline
\hline
Ejecutar / Probar & CmdR / CmdU \tabularnewline
Abrir la vista previa del canvas & OptionCmd \tabularnewline
Abrir un archivo al vuelo & CmdO \tabularnewline
Búsqueda global & CmdF \tabularnewline
Reformatear la sangría & CtrlI \tabularnewline
Comentar / Descomentar & Cmd/ \tabularnewline
Limpiar la caché de compilación & CmdK \tabularnewline
Ir a la definición & Cmd + clic \tabularnewline
Renombrar (cambio de nombre global y seguro) & Clic derecho Refactor Rename \tabularnewline
\hline
\end{tabularx}
\end{center}

\begin{center}\small\color{gray}Atajos de teclado de Xcode de uso frecuente\end{center}

\begin{tcbinfo}{Consejo: si la compilación se vuelve misteriosa, limpia primero la caché}
Los fenómenos paranormales del tipo «ayer funcionaba y hoy ya no» se deben, seis o siete de cada diez veces, a una caché de compilación corrompida. Con CmdK la limpias, vuelves a compilar y eso cura la mayoría de los misterios. Si no basta, reinicia Xcode (y de paso el Mac); si aún así nada, «borra la carpeta DerivedData»: es el trío definitivo que los desarrolladores iOS se pasan de boca en boca.

\end{tcbinfo}

\booksection{Recursos de aprendizaje recomendados}

\booksection{Recursos de aprendizaje recomendados}

\begin{center}
\begin{tabularx}{\linewidth}{|X|X|X|}
\hline
\textbf{Tipo} & \textbf{Recurso} & \textbf{Por qué lo recomendamos} \tabularnewline
\hline
Oficial & developer.apple.com & Documentación oficial de Apple; contiene los tutoriales de SwiftUI/Swift, la máxima autoridad \tabularnewline
Oficial & SwiftUI Tutorials de Apple & Curso oficial que te guía paso a paso, amabilidad máxima con el principiante \tabularnewline
Oficial & Sitio oficial de Swift (swift.org) & Documentación oficial del lenguaje y playground en línea \tabularnewline
Comunidad & Hacking with Swift & Biblioteca de tutoriales gratuitos de Paul Hudson, calidad de primera fila reconocida por todos \tabularnewline
Comunidad & Stack Overflow (etiqueta swiftui) & La primera parada cuando un error te deja sin salida \tabularnewline
Vídeo & Canal oficial de Apple Developer (vídeos de la WWDC) & Explicaciones autorizadas de las novedades de cada año \tabularnewline
Diseño & Guía de interfaz humana de Apple (HIG) & La biblia oficial para lograr interfaces con «sabor Apple» \tabularnewline
Complemento & Séptima parte, «Programación con IA en la práctica» (capítulos 37-47) & El método completo para desarrollar iOS con ayuda de la IA \tabularnewline
\hline
\end{tabularx}
\end{center}

\begin{center}\small\color{gray}Lista de recursos para aprender desarrollo iOS (agosto de 2026)\end{center}

\booksection{Palabras finales: de la primera app a la app número N}

Mirando atrás: pasaste de «no saber qué es SwiftUI» a instalar Xcode, aprender Swift, dibujar interfaces con SwiftUI, conectar tu app a una base de datos y a la red, y después probarla, empaquetarla y publicarla — en este camino ya recorriste el tramo más difícil del paso de cero a uno.

El mundo de la tecnología cambia muy rápido: los frameworks se actualizan, las herramientas se renuevan y la IA irá asumiendo cada vez más trabajo de «teclear». Pero hay algo que no caduca: tu capacidad de descomponer una idea difusa en pasos claros y materializarla poco a poco. Las herramientas son el amplificador; el razonamiento, el cimiento. El amplificador se puede sustituir; el cimiento no debe tambalearse.

Ya no eres solo "alguien que sabe usar el móvil": puedes convertir las ideas de tu cabeza en apps que funcionan de verdad en cientos de millones de iPhone. Desde hoy, crea con el teclado: ese es el romanticismo de los desarrolladores iOS.

\booksection{Referencias}

\begin{enumerate}
  \item Apple Inc. (2026). SwiftUI Framework Documentation. developer.apple.com.
  \item Apple Inc. (2026). Swift Programming Language Guide. swift.org.
  \item Apple Inc. (2026). SwiftData Framework Documentation. developer.apple.com.
  \item Apple Inc. (2026). App Store Review Guidelines. developer.apple.com/app-store/review.
  \item Apple Inc. (2026). Human Interface Guidelines (HIG). developer.apple.com/design.
  \item Apple Inc. (2026). Xcode 27 Release Notes and App Distribution Guide.
\end{enumerate}
